% TOP_PROBLEM: {"problemId": "problem.birch-and-swinnerton-dyer-conjecture", "canonicalTitle": "Birch and Swinnerton-Dyer Conjecture", "releaseRank": 6, "primaryDomain": "number_theory_arithmetic_geometry", "primaryDomainLabel": "Number theory & arithmetic geometry", "displayStatus": "open", "releaseStatus": "open"}
% !TeX program = lualatex
% Definition and sources only; this file is not an LLM solution attempt.
\documentclass[11pt]{article}
\usepackage[margin=25mm]{geometry}
\usepackage{fontspec}
\setmainfont{Latin Modern Roman}
\usepackage{amsmath,amssymb,mathtools}
\usepackage{unicode-math}
\setmathfont{Latin Modern Math}
\usepackage{xurl,hyperref}
\hypersetup{colorlinks=true,urlcolor=blue}
\setcounter{secnumdepth}{0}
\setlength{\parindent}{0pt}
\setlength{\parskip}{5pt}
\setlength{\emergencystretch}{3em}
\begin{document}
\section*{\texorpdfstring{Birch and Swinnerton-Dyer Conjecture}{Birch and Swinnerton-Dyer Conjecture}}\label{6}
\subsection{Definitions and mathematical statement}
An elliptic curve $E/{\mathbb{Q}}$ is a smooth projective genus-one curve with a specified rational point. Its rational points form a group $E({\mathbb{Q}})\simeq{\mathbb{Z}}^r\oplus E({\mathbb{Q}})_{\rm tors}$; $r$ is its Mordell--Weil rank. At a good prime $p$, put $a_p=p+1-\#E({\mathbb{F}}_p)$ and use the factor $(1-a_pp^{-s}+p^{1-2s})^{-1}$ in $L(E,s)$. The usual bad-prime factors complete the Hasse--Weil Euler product. Write $r_{\rm an}=\operatorname{ord}_{s=1}L(E,s)$, the smallest derivative order with a nonzero value at one. The selected assertion is
\[
 \forall E/{\mathbb{Q}},\qquad \operatorname{rank}E({\mathbb{Q}})=\operatorname{ord}_{s=1}L(E,s).
\]
It does not include the stronger leading-coefficient formula.
\par\smallskip\textit{Formulation and concept references:} \href{https://www.claymath.org/library/monographs/MPPc.pdf}{[S1]}, \href{https://www.claymath.org/wp-content/uploads/2022/05/birchswin.pdf}{[E1]}.

\subsection{Short English statement}
Does the number of independent rational points on an elliptic curve equal the multiplicity of its $L$-function's zero at one?
\subsection{Sources}
\begin{itemize}
\item[S1]Clay Mathematics Institute, The Millennium Prize Problems.
\url{https://www.claymath.org/library/monographs/MPPc.pdf}.
\textit{Formulation source.}
\item[S2]Birch and Swinnerton-Dyer Conjecture.
\url{https://www.claymath.org/millennium/birch-and-swinnerton-dyer-conjecture/}.
\textit{Further reference; not a proof endorsement. Consulted 24 September 2026: Institutional overview only.}
\item[S3]Introduction to the Birch and Swinnerton-Dyer Conjecture — Fabio Ferrari Ruffino.
\url{https://arxiv.org/abs/2608.25975}.
\textit{Further reference; not a proof endorsement.}
\item[S4]Proof of the Birch and Swinnerton-Dyer Conjecture — Yoshinori Shimizu.
\url{https://zenodo.org/records/17010555}.
\textit{Further reference; not a proof endorsement.}
\item[E1]Andrew Wiles, The Birch and Swinnerton-Dyer Conjecture, official Clay Mathematics Institute problem description, pages 1–2.
\url{https://www.claymath.org/wp-content/uploads/2022/05/birchswin.pdf}.
\textit{Added primary source: Elliptic curves, the Hasse–Weil L-function, and the rank formulation; the leading-coefficient refinement is not added to this entry.. Consulted 24 September 2026.}
\end{itemize}
\end{document}
